% ECONOMICS_PROBLEM: {"id": "L11-337", "title": "Concentration versus market power", "jel_code": "L11", "source_id": "OP-0390"}
% !TeX program = xelatex
\documentclass[11pt,a4paper]{article}
\usepackage[margin=23mm]{geometry}
\usepackage{fontspec}
\setmainfont{Latin Modern Roman}
\usepackage{amsmath,amssymb,mathtools}
\usepackage{xcolor,enumitem,booktabs,longtable,array,xurl,hyperref}
\definecolor{muted}{HTML}{53636C}
\hypersetup{colorlinks=true,urlcolor=blue,linkcolor=blue}
\setlength{\parindent}{0pt}
\setlength{\parskip}{5pt}
\newcommand{\R}{\mathbb R}
\newcommand{\E}{\mathbb E}
\newcommand{\N}{\mathbb N}
\newcommand{\ind}{\mathbf 1}
\newcommand{\Var}{\operatorname{Var}}
\newcommand{\Cov}{\operatorname{Cov}}
\newcommand{\supp}{\operatorname{supp}}
\DeclareMathOperator*{\argmin}{arg\,min}
\DeclareMathOperator*{\argmax}{arg\,max}
\newcommand{\entrymeta}[3]{{\sffamily\small\color{muted}\textbf{\texttt{#1}}\quad|\quad #2\par #3\par}}
\newcommand{\Problemref}[1]{\texttt{#1}}
\newcommand{\JELref}[2]{\texttt{#2} (JEL \texttt{#1})}
\begin{document}
\section*{Concentration versus market power}
\textbf{Problem ID:} \texttt{L11-337}\quad \textbf{JEL:} \texttt{L11}\par
% BEGIN ECONOMICS STATEMENT
\label{problem:OP-0390}
\label{audited:ECON-033}
{\sffamily\small\color{muted}Original subject group: Industrial organization and competition\par}
\label{definitions:problem:30007672}
\textbf{Audit classification:} Published research agenda. The original two-part decomposition cannot exhaust an observed change when demand, population or product quality also change. A welfare functional and counterfactual equilibrium selection were missing.
\subsection*{Definitions and precise research target}
Consider United States retail-product markets \(m=1,\ldots,M\) observed annually, where \(M\) is the number of sampled markets. Define each market by substitutable products and a fixed geographic boundary. Let \(s_{jmt}\) be firm \(j\)'s expenditure share in market \(m\) at year \(t\), and let \(H_{mt}=\sum_j s_{jmt}^{2}\) measure concentration. Observe transaction prices, quantities, product characteristics, ownership and input-cost shifters; marginal costs and competitive conduct remain latent. A structural parameter \(\theta\) determines demand, costs and conduct, and \(P_\theta\) denotes its implied distribution of these observations. For an observed distribution \(P\), define \(\Theta(P)=\{\theta:P_\theta=P\}\). Construct a transparent decomposition of the cost-and-conduct change in consumer surplus holding demand and the consumer population fixed into a technology component, obtained by changing costs with conduct fixed, and a conduct component, obtained by changing conduct with costs fixed. Consumer surplus is expected money-valued utility net of expenditure in the defined market. Specify the order of these counterfactual changes because interactions make decompositions path dependent. The task is to determine which additional instruments or restrictions make the sign of the conduct component invariant across \(\Theta(P)\), and otherwise report its identified range. Rising \(H_{mt}\) alone is an observation rather than a maintained causal explanation.

\textbf{Definitions, assumptions and mathematical qualifications.} A market definition is an input: a product set, geographic set and outside option fixed before observing outcomes. Firm shares satisfy \(s_{jmt}\geq0\) and \(\sum_j s_{jmt}=1\). Specify quasi-linear utility \(v_i(x)-p\cdot x\), population weights and an equilibrium-selection kernel; consumer surplus is the expectation of indirect utility relative to the outside option. Write \(CS(c,k;d)\) for surplus under cost technology \(c\), conduct \(k\), and demand/population \(d\). For a fixed reference \(d\), define \(T=CS(c_1,k_0;d)-CS(c_0,k_0;d)\) and \(K=CS(c_1,k_1;d)-CS(c_1,k_0;d)\). Then \(T+K\) equals the cost-and-conduct contrast. Changing demand additionally requires a third component. All counterfactual equilibria must exist and the nonempty model class must reproduce the observed law. Its conduct range is \([\inf_{\theta\in\Theta(P)}K_\theta,\sup_{\theta\in\Theta(P)}K_\theta]\).
\subsection*{Variants and assumption changes}
\begin{enumerate}
\item \textbf{Editorial, fixed demand.} Compare the two possible orders of cost and conduct changes; average their two conduct contributions to obtain a two-factor Shapley allocation. Test sign invariance separately for each ordering.
\item \textbf{Editorial, changing demand.} Introduce \(d_0,d_1\) as a third factor, average contributions over all six orders, and identify the technology, conduct and demand components jointly.
\end{enumerate}
\subsection*{References and source audit}
\textbf{Checked source.} 24 August 2024 author draft, Section 6, printed pp. 35--37 (market boundaries, technology and conduct). Market-specific conduct, markup validation and merger research are supported. New mathematical targets are editorial. Full primary text retrieved. \url{https://web.stanford.edu/~ayurukog/trendsincompetition.pdf}.
\begin{enumerate}\raggedright
\item\hypertarget{definitions:source:30007672:1}{} Carl Shapiro; Ali Yurukoglu. 2024. \emph{Trends in Competition in the United States: What Does the Evidence Show?}. 24 August 2024 author draft, Section 6, printed pp. 35--37 (market boundaries, technology and conduct).
\url{https://web.stanford.edu/~ayurukog/trendsincompetition.pdf}.
DOI: \url{https://doi.org/10.3386/w32762}.
\end{enumerate}

\subsection*{Integrated refinement: ECON-033}
{\sffamily\small\color{muted}Efficiency versus market-power explanations of concentration\par}
This refinement adopts the additional source conventions
(page~\pageref{audited:conventions}), subject to its local restrictions.
\textbf{Concentration and welfare decompositions on the same counterfactual path.}
For productivity $a$, conduct $\kappa$, equilibrium shares $s_j(a,\kappa)$ and
$H(a,\kappa)=\sum_js_j(a,\kappa)^2$, also identify
\[
\Delta H_a=H(a_1,\kappa_0)-H(a_0,\kappa_0),\qquad
\Delta H_\kappa=H(a_1,\kappa_1)-H(a_1,\kappa_0).
\]
Bound both components and the associated consumer-welfare effects using the same
model class, population and equilibrium selector as the existing surplus
decomposition. Fix demand and market definitions or introduce demand explicitly
as an additional factor. Components depend on ordering; reverse the order or use
a declared averaging rule when assessing robustness. Concentration alone does
not identify either efficiency or conduct changes.
\subsection*{Supplied source audit for ECON-033}
Local [S1], [S2], etc. in this audit and the references immediately below
belong to ECON-033, independently of the original entry's reference numbering.
[S1] and [S2], Sections II.A and IV, support distinguishing concentration from market power. The path-specific decomposition is editorial.
\subsection*{References for integrated refinement ECON-033}
{\footnotesize\raggedright
\par\noindent\textbf{[S1]} Chad Syverson (2019). \emph{Macroeconomics and Market Power: Context, Implications, and Open Questions}. Journal of Economic Perspectives 33(3), 23–43. DOI: 10.1257/jep.33.3.23.\par \textit{Verified locator / source role:} Discussion of concentration, market power, and alternative explanations.\par \url{https://www.aeaweb.org/articles?id=10.1257/jep.33.3.23}\par\smallskip
\par\noindent\textbf{[S2]} Chad Syverson (2019). \emph{Macroeconomics and Market Power: Facts, Potential Explanations, and Open Questions}. Brookings Economic Studies, January 2019; author research essay.\par \textit{Verified locator / source role:} Section II.A, printed pp. 4–8; Section IV: concentration versus competition.\par \url{https://www.brookings.edu/wp-content/uploads/2019/01/ES_20190116_Syverson-Macro-Micro-Market-Power.pdf}\par\smallskip
}

\subsection*{Common conventions: Source mathematical conventions}

These conventions apply to each entry unless explicitly replaced.

\textbf{Models and identification.} A primitive model \(m\) specifies all probability laws, feasible choices, information, equilibrium and measurement rules used by its target. The admissible class \(\mathcal M\) is fixed independently of outcomes. If \(P_O(m)\) is its observable probability law and \(\tau(m)\) the target, its identified set at observable law \(P\) is
\[
 \mathcal I_\tau(P)=\{\tau(m):m\in\mathcal M,\ P_O(m)=P\}.
\]
Point identification means that this set is a singleton. An empty set signals incompatibility of data and assumptions. Coordinate bounds need not describe the joint set. Bounds are sharp only if every retained target is attainable or arbitrarily approachable by compatible models. Finite bounds require explicit restrictions on unobserved quantities; unrestricted confounding, hidden wealth, extrapolation or arbitrary equilibrium selection can yield uninformative sets.

\textbf{Causal interventions.} A potential outcome \(Y_i(p)\) is the outcome under a complete policy regime \(p\), including timing, financing, information and market spillovers. The target population and units are fixed before treatment. With interference, use the complete assignment vector or a justified exposure mapping; individual no-interference notation alone cannot identify an economy-wide intervention. A difference of mean potential outcomes does not require the cross-world joint distribution. Conditioning on post-treatment migration, survival, enrollment or employment changes the estimand and needs separate justification.

\textbf{Statistical statements.} An estimator, confidence set, forecast or test specifies a sampling law, model class and loss/coverage criterion. Uniform \((1-\alpha)\)-coverage means \(\inf_{m\in\mathcal M}\Pr_m\{\tau(m)\in C_n\}\geq1-\alpha\), or an explicitly asymptotic version. The whole parameter space trivially has coverage, so informativeness requires a stated width, risk or power objective. A forecasting improvement is relative to a named benchmark, information set and evaluation population.

\textbf{Equilibrium and optimization.} An equilibrium comprises feasible plans, optimality, beliefs where needed, budget constraints and all market/resource clearing. The primitives or a stated benchmark define each correspondence \(\mathcal E(p;m)\). Nonemptiness is assumed or proved, never inferred just from compact policy sets. A selected-equilibrium target uses a declared selection rule; a robust target quantifies over all permitted equilibria. Use suprema or infima unless attainment is established. An \(\varepsilon\)-optimal policy is feasible and within \(\varepsilon\) of the finite optimum value. Report an empty feasible set separately.

\textbf{Welfare and accounting.} Normative utility, interpersonal normalization, social weights, numeraire, discounting and the use of public receipts are primitives. Behavioral utility need not coincide with the welfare criterion. Household welfare includes ownership income and rents once; adding profits again to compensating variation generally double-counts them. A monetary equivalent is defined by an explicit transfer and price convention, with monotonicity and range conditions. Average effects, mechanism contributions, transfers and welfare losses are different targets. Infinite sums require convergence and dynamic feasible plans require terminal or no-Ponzi conditions.

\textbf{Source versions.} A working-paper date, revision date and journal publication year are distinguished. A DOI for a journal version is not automatically the DOI of an earlier manuscript. Locators refer to the stated version and printed pages where specified. A metadata-only check does not verify a claim about a full-text conclusion. Every entry's source subsection narrows the provenance claim accordingly.

\subsection*{Common conventions: Additional-source status and mathematical conventions}

\label{audited:conventions}

Of the 100 supplied ECON entries, 65 distinct problems appear here; 32 close
matches refine earlier entries, and three duplicate questions map to existing
entries. Appendix~\ref{app:audited-crosswalk} lists every destination.
The attachment's P/R/S/U distinctions and source audits are retained. Its
precise mathematical formalizations are editorial unless the individual audit
says otherwise. Candidate subsidy targets ECON-004--006 retain their U flags;
the resolved approval-core question ECON-007 updates OP-0202 above.
The following supplied conventions govern both this part and the attributed
ECON refinements elsewhere in the catalogue, subject to local restrictions.

\textbf{Parameterized questions.} A problem family lists inputs that must be fixed before an instance is evaluated: population, horizon, policies, information, data law, model class and welfare weights. These are required choices, not quantities the citations are assumed to specify. Undefined applications should not be treated as identified estimands.

\textbf{Indices and integrability.} Sets of agents, firms and regions are finite unless a population measure is explicitly used. \(i,j\) index units, \(r\) regions, \(t\) dates and \(h\) horizons. \(\mathbb E\) and \(\Pr\) refer to the declared population and law; \(\mathbf1\{A\}\) is an indicator. Every displayed expectation and discounted sum needs existence in the stated domain. Infinite horizons use a discount in \((0,1)\) and a convergence condition; finite horizons may use discount one.

\textbf{Optimization and existence.} A displayed \(\max\), \(\min\), or \(\arg\max\) is conditional on attainment. For finite feasible menus this is immediate; general spaces need continuity and compactness/coercivity or another existence argument. Without attainment, the target is the supremum/infimum and arbitrarily near-optimal feasible choices. \(\inf\varnothing=+\infty\). Empty feasibility and an unidentified target are different issues.

\textbf{Potential outcomes.} \(Y_i(z)\) is the outcome under a specified intervention, not the observed conditional mean among people choosing \(z\). In binary comparisons \(\Delta Y_i=Y_i(1)-Y_i(0)\). Specify assignment, treatment versions, compliance, information and observation horizon. Interference is allowed under a full policy vector; compressed exposure notation needs a justified mapping. No interference, consistency and overlap are substantive assumptions, not automatic consequences of notation.

\textbf{Populations and observation.} Fix baseline eligibility and population weights. Conditioning on post-policy employment, survival, relocation or program uptake can change the population being compared. Death is not ordinary loss to follow-up; outcomes truncated by death need a composite convention, joint survival target, or explicitly defined survivor stratum. Fertility-changing interventions need a population functional rather than an unexplained fixed descendant cohort. Measurement and missingness models must be supplied.

\textbf{Identification.} A structural model \(m\in\mathcal M\) implies observable law \(P_m\) and target \(\theta(m)\). Identification means
\[
 P_{m_1}=P_{m_2}\ \Longrightarrow\ \theta(m_1)=\theta(m_2)
 \quad\text{for all }m_1,m_2\in\mathcal M .
\]
Without identification, the target is
\[
 \Theta(P;\mathcal M)=\{\theta(m):m\in\mathcal M,\ P_m=P\}.
\]
Writing \(\theta(P)\) before establishing identification can hide incompatible latent models. Confidence sets additionally need a sampling law, confidence level, asymptotic sequence and uniformity class. Point identification, coverage and informative precision are separate properties.

\textbf{Mediation.} Total effects of randomized assignment are distinct from cross-world natural indirect effects. Belief/effort-channel effects require a meaningful mediator intervention and additional measurement, exclusion or mediation assumptions. Randomization of the initial policy alone does not supply those assumptions. Report bounds or interventional alternatives when the stronger target is unsupported.

\textbf{Equilibrium.} A counterfactual \(W(z)\), \(p(z)\), or \(\mu_t^z\) requires individual optimization, feasibility, government/resource budgets, market clearing and state-transition laws. State equilibrium existence and selection, or report ranges over compatible equilibria. Initial-state integration includes subsequent endogenous transitions; current-state integration must use the policy-dependent distribution. Reduced-form invariance after policy is not assumed.

\textbf{Welfare accounts.} Scores, utility, health, dollars and ecological quantities require explicit conversion before addition. Fees, taxes, subsidies, wages and extortion payments are transfers between parties; distributional weights can make them welfare relevant, but their face value is not automatically a resource loss. State ownership, geographic boundaries, foreign-income weights and fiscal finance. Do not add profits to household welfare already including dividends, or count land capitalization and the underlying benefit twice. A benefit-cost expression must distinguish gross benefits from net welfare and subtract every real cost exactly once.

\textbf{Derivatives and risk.} Log derivatives require positive arguments; local responses require admissible perturbations and specified equilibrium branches. Differentiation under expectations needs regularity/domination, and boundaries may allow only one-sided derivatives. Risk uses a specified law; ambiguity uses a declared nonempty law set on a common history space. Adaptive policies must be nonanticipative. A robust criterion is an editorial choice unless attributed explicitly.

\textbf{Scope overlaps.} Allocation, collective choice, contracts, household bargaining and coordinated policy overlap economics with optimization and game theory. They are retained because they appeared in the original catalogue; no claim of a strictly game-theory-free selection is made.
% END ECONOMICS STATEMENT
\end{document}
